\chapter{Prototypische Umsetzung}
\label{Implementierung}
% GEÄNDERT: Vorlagentext entfernt; Kapitel spiegelt die Teilaspekte aus Kap. 3 (Betreueranmerkung S4: gleiche Überschriften, keine Überraschungen, nur was die Evaluation benötigt)

% NEU: Kapiteleinleitung
Aufbauend auf dem in Kapitel~\ref{Ansatz} beschriebenen Zielentwurf beschreibt dieses Kapitel dessen prototypische Umsetzung. Die Implementierung dient dabei ausschließlich der Evaluation in Kapitel~\ref{Evaluation}, sodass nur jene technischen Details dargelegt werden, die zum Nachvollziehen der Ergebnisse sowie zum Nachbau des Verfahrens erforderlich sind. Die Gliederung folgt den Teilaspekten aus Abschnitt~\ref{sec:teilaspekte}, und für jeden Teilaspekt wird ausgewiesen, welche Bestandteile des Zielentwurfs umgesetzt sind und welche nicht.

\section{Systemumgebung und Datenbasis}
\label{sec:impl-umgebung}

% NEU (übernommen aus Textbaustein 21.09.2026, Abschnitt 4.1)
Die prototypische Umsetzung erfolgt gegen ein SAP-System, auf das nicht unmittelbar, sondern über eine in der Entwicklungsumgebung hinterlegte Verbindungsbeschreibung zugegriffen wird. Der Zugriff auf die Metadatenebene erfolgt über die REST-Schnittstelle der ABAP Development Tools (ADT); die fachlichen Abfragen werden über deren Endpunkt für freie Datenvorschauen abgesetzt.

Die Datenbasis der semantischen Suche bildet das Data Dictionary. Herangezogen werden die Feldliste der Tabellen sowie die Beschreibungstexte auf Tabellen- und auf Datenelementebene, sodass jedes Feld mit den auf drei Ebenen verfügbaren sprachlichen Informationen beschrieben wird. Die Eingrenzung des betrachteten Objektbestands erfolgt über den Objektkatalog, der jedem Entwicklungsobjekt sein Paket sowie seinen Ursprung zuordnet. Diese Eingrenzung ist von der in Abschnitt~\ref{subsec:ta-erschliessen} begründeten inhaltlichen Filterung des Schemakorpus zu unterscheiden: Der Objektkatalog begrenzt, welcher Ausschnitt des Systems überhaupt betrachtet wird, während die inhaltliche Filterung innerhalb dieses Ausschnitts bestimmt, welche Tabellen als Träger von Customizing in den Korpus aufgenommen werden.
% Zu prüfen vor Übernahme: Nennung von Systembezeichnung und Mandant mit dem Kooperationspartner klären.
% Ergänzen: Ablage der Codebasis und Folgen für die Nachvollziehbarkeit (Textbaustein 15.09.2026).

\section{Relevante Konfiguration erschließen}
\label{sec:impl-erschliessen}
% TODO (Umsetzung, Stand Prototyp 23.09.2026):
%  - Korpusaufbau offline: Pakete über TADIR -> Filter DD02L (TABCLASS = TRANSP, CONTFLAG C/G) -> Felder DD03L/DD08L/DD02T -> Texte DD04T mit Sprachrückfall -> F1-Doku DOKTL
%  - Granularität Feld (Default) bzw. Tabelle; Baseline ohne F1-Doku
%  - Einbettung über Orchestration-Dienst des Generative AI Hub; Dimension konfigurierbar; L2-Normalisierung
%  - FAISS IndexFlatIP (exakte Suche, Skalarprodukt = Kosinus) + Metadaten-JSON
%  - Werkzeuge search_schema, describe_table, query_table (nur Gleichheitsfilter, Feldprojektion, beleg_sql)
%  - MCP-Server (FastMCP), streamable HTTP, einmal gestartet je Sitzung
%  - Umfang des Korpus (Anzahl Tabellen/Felder) angeben

\section{Semantisch interpretieren}
\label{sec:impl-interpretieren}
% TODO (Umsetzung):
%  - ReAct-Agent (LangChain create_agent), Router customizing | code, Systemprompt mit Vorgehensschritten
%  - Modellanbindung über Generative AI Hub, gekapselt in einer Funktion (Modellwechsel)
%  - Klartextauflösung: CHECKTABLE -> Texttabelle (FRKART TEXT), Join nur bei eindeutiger Zuordnung und Wertüberlappung
%  - Merkmale: is_feature (T549B), decode_feature (T549C), resolve_value text | lexical | abstain
%  - read_source, search_objects, list_package für kundeneigenen Code
%  - Oberfläche (Chainlit) mit Protokoll und Belegen

\section{Aussagen an Systemdaten absichern}
\label{sec:impl-absichern}
% TODO (Umsetzung): Belegführung über beleg_sql und Trace-Sammlung; Prüfschleife als LangGraph-Zustandsautomat;
% Stand ehrlich ausweisen: im produktiven Agentenpfad derzeit Belegdisziplin über den Systemprompt, deterministische Schleifen als Vergleichsbasis.

\section{Abweichung vom Standard ausweisen}
\label{sec:impl-abweichung}
% TODO (Umsetzung): derzeit nur Namensraum (Z/Y-Datenelemente, Z-Objektsuche); Referenzabgleich noch nicht umgesetzt.
